\chapter*{LỜI MỞ ĐẦU}\addcontentsline{toc}{chapter}{Lời mở đầu}
Chọn trường và chọn ngành là hai quyết định liên kết nhưng không hoàn toàn đồng nhất. Trường đại học tạo ra môi trường, nguồn lực và tín hiệu danh tiếng; ngành học định hình nội dung chuyên môn, trải nghiệm học tập và cơ hội nghề nghiệp. Với học sinh lớp 12, quyết định thường được hình thành trong điều kiện thông tin chưa đầy đủ, chịu ràng buộc bởi năng lực, tài chính gia đình, khoảng cách địa lý và kỳ vọng của người thân.

ĐBSCL có đặc điểm địa bàn rộng, chi phí di chuyển đáng kể giữa các tỉnh, cơ cấu kinh tế đang chuyển đổi và chênh lệch cơ hội tiếp cận thông tin. Vì vậy, một mô hình đồng thời xem xét cả chọn trường và chọn ngành có giá trị cho công tác hướng nghiệp và tuyển sinh. Luận văn được tổ chức thành ba chương: Chương 1 trình bày vấn đề, mục tiêu và thiết kế nghiên cứu; Chương 2 xây dựng cơ sở lý thuyết, thang đo và giả thuyết; Chương 3 báo cáo quy trình phân tích, kết quả và thảo luận. Phần cuối tổng hợp kết luận, hàm ý, giới hạn và hướng nghiên cứu tiếp theo.
